This chapter describes the prototype implementation of the designed system, as well as its testing and evaluation.

For the entire implementation, TypeScript, Node.js, and npm were used. The reason behind this choice is both prior experience with this ecosystem of the author of this thesis, as well as its popularity for full-stack web development, ensuring high amount of useful supporting resources.

TypeScript~\cite{typescript} is a popular \enquote{strongly typed programming language that builds on JavaScript.} It \enquote{adds additional syntax to JavaScript to support a tighter integration with (...) an editor} to \enquote{catch errors early.} The code \enquote{converts to JavaScript, which runs anywhere JavaScript runs: in a browser, on Node.js, Deno, Bun.} TypeScript is a static type checker, meaning that it performs type checking only at compile time, not during the runtime.

\begin{sloppypar}
Node.js~\cite{node_js} is a \enquote{cross-platform JavaScript runtime environment that lets developers create servers, web apps, command line tools and scripts.} It is \enquote{built on the V8 JavaScript engine} (which is also used~\cite{v8} by Chromium-based web browsers, such as Google Chrome). The Node.js documentation~\cite{node_js} further states that \enquote{npm is the standard package manager for Node.js,} with npm~\cite{npm} claiming that it is \enquote{the largest software registry in the world.}
\end{sloppypar}

Each application that is part of the implementation is kept in a separate directory: \textit{booking-service} for the booking service, \textit{booking-api-gateway} for the API gateway, and \textit{booking-client-app} for the client application. These directories with the source code are included in the content of the digital attachment~\ref{part:contents_of_the_digital_attachment}.

\begin{sloppypar}
Every application directory contains a \textit{package.json} file and an \textit{npm-shrinkwrap.json} file. These files in the JavaScript Object Notation (JSON) data format are both related to npm, with \textit{package.json}~\cite{npm_docs} describing many things, including general metadata about the project (package name, version adhering to semantic versioning, license information, author information, important URLs, etc.), npm scripts' definitions, lists of npm dependencies (of various types) along with their version, application entrypoints, and information about supported environments. The file \textit{npm-shrinkwrap.json} is then generated by npm when installing dependencies (using the \mintinline{sh}{npm install} command) and it is a publishable variant of the \textit{package-lock.json} file used by default to describe \enquote{the exact (dependency) tree that was generated, such that subsequent installs are able to generate identical trees, regardless of intermediate dependency updates.} This ensures consistently reproducible dependency installations. Dependencies are installed into the \textit{node\_modules} directory of each project (this directory is not included in the source code).
\end{sloppypar}

To fulfill non-functional requirement~\ref{req:open_source}, each application directory also contains a \textit{LICENSE.txt} file with the popular permissive open-source MIT license (which is also correctly denoted in the \textit{package.json} files) and each application only uses dependencies with compatible permissive open-source licenses.

Furthermore, the file \textit{Containerfile} is included in each application directory. \textit{Containerfile}~\cite{podman_github} (which is an alternative name for \textit{Dockerfile}) is \enquote{a configuration file that automates the steps of creating a container image.} Application container images can then be created using tools like Docker, Podman, and Buildah, either in the Docker image format or in the Open Container Initiative (OCI) image format (which is largely similar). These container images can then used by orchestration tools (such as Kubernetes) for~\cite{kubernetes} \enquote{automating deployment, scaling, and management of (...) applications}. Additionally, the \textit{compose.prod.yaml} Compose file is attached to enable easily building and running all the service containers in production mode using tools like Docker Compose, and the \textit{booking-service} directory contains a \textit{compose.yaml} file with the dependant services for the booking service to run (and thus is useful for development). This fulfills non-functional requirement~\ref{req:cloud_native}.

Lastly, each application directory contains a \textit{README.md} file in the Markdown format, describing how to run it, and a \textit{.env.template} file with a template for the environment variables that need to be set for the application to run.

\section{Booking Service}
\label{part:i_booking_service}

The booking service was built around the framework NestJS (Nest). Nest~\cite{nestjs} \enquote{is built with and fully supports TypeScript} and is used \enquote{for building (...) Node.js server-side applications.} \enquote{Under the hood, Nest makes use of robust HTTP server frameworks like Express (the default) and optionally can be configured to use Fastify as well.} As Express is the default and (judging by download numbers from npm) the much more popular choice, and the Nest documentation~\cite{nestjs} occasionally mentions unsupported functionality when using Fastify, Nest was configured to use Express than the alternative. Nest \enquote{provides an out-of-the-box application architecture which allows developers and teams to create highly testable, scalable, loosely coupled, and easily maintainable applications.}

The framework makes frequent use of the dependency injection pattern. This pattern is used to inject services into controllers, which are used to handle incoming requests. The controllers are responsible for handling incoming requests and returning responses. The services are responsible for handling the business logic of the application. The services are also responsible for handling the communication with the database.

For Object Relational Management (ORM), MikroORM was used. For the database, PostgreSQL was used, which is an ACID compliant relational database management server, with support for JSON data types. The diagram of the PostgreSQL database can be seen in the figure~\ref{fig:db_diagram}.

\begin{figure}
    \centering
    \includegraphics[width=1.0\textwidth]{content/implementation/db_diagram.uxf.svg.pdf}
    \caption[Diagram of the booking service database]{Diagram of the booking service database}
    \label{fig:db_diagram}
\end{figure}

The booking service implements the booking protocol described in~\ref{part:booking_protocol}. Furthermore it offers internal APIs for use by the client.


\section{API Gateway}
\label{part:i_api_gateway}

The API gateway behaves to the outside worlds as a booking service. It lets anybody access the booking protocol APIs of the booking service, but only authenticated users to use to internal APIs meant for the booking client app.

Much like the booking service, the API gateway was built around the framework NestJS (Nest), with Express enabled.

For user authentication, the library SuperTokens was used. For proxying requests to booking service, the library http-proxy-middleware was used.

The API gateway also provides on internal endpoint \texttt{POST /forward-booking-request} for use by the booking service only, which is used to forward booking requests from the booking service to remote booking services.


\section{Client Application}
\label{part:i_client_application}

The booking client application shall use a single page application (SPA) architecture. The SPA architecture is a design pattern that allows the client application to load a single HTML page and dynamically update the page as the user interacts with the application. This design pattern allows the client application to be more responsive and provide a better user experience.

The client application will be a progressive web application (PWA) that can be installed on the user's device and possibly even work offline (though that will not be implemented as of now). This will also enable the use of the custom booking URL schema (on supported platforms). And as any web application, it will be accessible on many different devices and platforms.

Lastly, it will be a mobile-first design to ensure that the application is usable on mobile devices.

Initially, the Figma design tool was used to create wireframes of the UI, however, the UI has changed by a not significant amount since then, so the original design is left out of the thesis.


\section{Testing and Evaluation}
\label{part:testing_and_evaluation}

First, every part of the implementation was tested manually by the author of this thesis. For the client application, this meant clicking through every page and state, testing edge-case inputs (invalid usernames, booking addresses, form data) and even setting up booking service failure scenarios by temporarily modifying the booking service to respond with invalid data. For the API gateway and the booking service, the command line utility cURL was used to send both valid and invalid requests, including different HTTP header combinations.

Afterwards, when the implementation appeared stable, user usability testing was performed. 

\subsection{Usability testing}

Usability testing was performed with three subjects. \textbf{Subject A} is a 26 year old male software engineer with a university degree, who lives in the suburbs of Prague. \textbf{Subject B} is a 31 year old male army soldier, who lives in Prague. \textbf{Subject C} is a 52 year old female business manager, who lives in the suburbs of Prague and is not too technologically savvy.

The following scenario was chosen based on the designed use case and given to testing subjects: \enquote{You are a manager at a mid-sized business (a couple hundred employees). Your company wants to organize a conference and needs to keep track of how many and who can partake. You know that you yourself will be at the conference.}

Then, the following tasks were given to each subject:
\begin{enumerate}
    \item \enquote{Create a new account and booking address.}
    \item \enquote{Create a new inventory with basic information about yourself as a manager.}
    \item \enquote{Create a new inventory item for the conference.}
    \item \enquote{Make a new booking of the conference item, since you know that you will be participating.}
    \item \enquote{Check if you can find your booking in the inventory.}
\end{enumerate}

The testing was done in person, on a Windows 11 laptop that was brought to the testing setting, with all the required services set up and the booking client application open in the latest Google Chrome browser.

The testing subjects performed the tasks as follows:
\begin{enumerate}
    \item Account and address creation:
    \begin{itemize}
        \item \textbf{Subject A} signed in using their GitHub account, input a valid username and the first try and created their booking address.
        \item \textbf{Subject B} took a moment to find out how to switch the sign-in component to the sign up state, then created an account with email and password. They input a username and created their booking address, though at first try they put in an invalid character for the username, which was quickly corrected thanks to the instant UI feedback.
        \item \textbf{Subject C} tried to fill in the email and password for their new account into the sign in form. When that attempt was unsuccessful, they realized the presence of the button to switch the component to the sign up state and created their new account successfully. Then they put in a simple username and created a booking address without an issue.
    \end{itemize}
    \item Inventory creation:
    \begin{itemize}
        \item \textbf{Subject A} immediately spotted the burger icon button that opens the drawer with a link to the business section, where they transitioned, filled out the new inventory form and created their inventory.
        \item \textbf{Subject B} first opened the profile dialog, but their second choice was the burger icon button. The subject created their inventory without issue.
        \item \textbf{Subject C} took a little while to look around the page, but their first choice was the burger icon button. The rest of the process was flawless.
    \end{itemize}
    \item Item creation:
    \begin{itemize}
        \item \textbf{Subject A} navigated to the add item page and filled out the form. They were experimenting with the duration field, which they said has a too exotic format, though it was not needed for this task. Then they proceeded to create the item.
        \item \textbf{Subject B} navigated to the add item page, filled out the form and added the item to the inventory.
        \item \textbf{Subject C} also navigated to the add item page, filled out the form and added the item to the inventory without anything to mention.
    \end{itemize}
    \item Item booking:
    \begin{itemize}
        \item \textbf{Subject A} first clicked the newly created item, then realized that they need to navigate back to the home page and then to the new booking page. They at first did not realize, that they need to input the whole booking address, but thanks to the instant feedback of the UI, they managed to get to the items for booking page and complete the booking.
        \item \textbf{Subject B} experienced the same issues as \textbf{Subject A}, but managed to complete the booking.
        \item \textbf{Subject C} also struggled to get to the correct page and then to find out what to input as the booking address. Once they got over this step, the process was straightforward.
    \end{itemize}
    \item Booking overview:
    \begin{itemize}
        \item \textbf{Subject A} finished the task successfully without anything to mention.
        \item \textbf{Subject B} accessed the booking through their home page and thought that they completed the task. They had to be told, that this is not the correct place, otherwise they would not complete the task successfully.
        \item \textbf{Subject C} first opened the booking through their home page, then re-read the instructions and went correctly to their inventory to access the booking.
    \end{itemize}
\end{enumerate}

\subsection{Evaluation}

Through the conducted usability testing, some UI/UX blind spots were found. The issues some users faced with the sign in form could be attributed to the fact, that the form was in English and neither were native English speakers, so the two similar terms \enquote{sign in} and \enquote{sign up} can be confusing. Improvements in the client application could be made to better parse user values of certain metadata (enabling for example more human readable duration or opening hours format). Some sections of the client app are very similar, which can be also confusing. Likely the biggest hurdle is the exotic nature of booking addresses. This issue could be improved by displaying the users' booking addresses more prominently in the application.

These issues, however, are relatively minor and the implementation is fully functional and fulfills the designed use case.

